\documentclass{article} % say
\usepackage{amsmath}
\usepackage{gensymb}
\usepackage{hyperref}
\usepackage{tikz}
\usepackage{tkz-euclide}
\usetikzlibrary {babel,angles,quotes,positioning}
\usepackage[utf8]{inputenc}
\usepackage[russian]{babel}
\begin{document}

\textbf{Задача}: Дан треугольник $ABC$, причём $|AB| \neq |AC|$. В этом треугольнике $AL$~--- биссектриса, а $AM$~--- медиана. Докажите, что $|AL| < |AM|$. 

\clearpage

\textbf{Решение}: 

\begin{center}
\begin{tikzpicture}[scale=1.25]
\tkzDefPoints{0/0/A,3/-5/B,-3/-1/C}
\tkzDefMidPoint(B,C) \tkzGetPoint{M}
\tkzDefLine[bisector](C,A,B) \tkzGetPoint{l}
\tkzInterLL(A,l)(B,C) \tkzGetPoint{L}

\tkzDrawPolygon(A,B,C)
\tkzDrawLine[add=0 and 0](A,M)
\tkzDrawLine[add=0 and 0](A,L)

\tkzMarkAngles[mark=|,size=0.3](C,A,L)
\tkzMarkAngles[mark=|,size=0.5](L,A,B)

\tkzDrawPoints(A,B,C,M,L)
\tkzLabelPoints(B)
\tkzLabelPoints[right](A)
\tkzLabelPoints[below](C,L,M)
\end{tikzpicture}
\end{center}

Положим без потери общности, что $|AC| < |AB|$. Значит:

\begin{equation}
  \angle ACB > \angle ABC \label{neq1}
\end{equation}

Заметим, что $\angle ALC = \pi - \angle ACB - \angle LAC$, а $\angle ALM = \angle ALB = \pi - \angle LAB - \angle ABL = \pi - \angle LAB - \angle ABC$. Но так как $\angle LAC = \angle LAB$ ($AL$~--- биссектриса) и по неравенству \ref{neq1} выходит, что $\angle ALC < \angle ALM$.

$\angle ALC + \angle ALM = \pi$, а значит $\angle ALM$~--- тупой. Следовательно, в треугольнике $LAM$ сторона $AM$~--- наибольшая и, в частности, длинней $AL$. Что и требовалось доказать.

\textbf{Источник}: А.~П.~Киселёв, <<Элементарная геометрия. Книга для учителя>> (1980), страница 121, задача 244. \url{https://disk.yandex.ru/d/QqFzGPzxTVXGag}

\end{document}


